\label{sec:conclusion}

This paper has provided a comprehensive legal analysis of incumbency protection as a
redistricting criterion and its implications for algorithmic redistricting. The analysis
rests on three pillars of federal constitutional doctrine: \textit{Karcher v.\ Daggett}'s
identification of incumbency protection as a permissible but not mandatory objective;
\textit{Cox v.\ Larios}'s limitation of incumbency protection to contexts where it
justifies only minor population deviations; and \textit{Thornburg v.\ Gingles}'
establishment of VRA \S{}2 requirements that are independent of and take precedence
over incumbency protection in majority-minority district contexts.

From these doctrinal foundations, we derived the algorithmic indifference standard:
a redistricting methodology that receives no incumbency data, optimises only constitutionally
required or permitted criteria, and produces incumbency outcomes consistent with the
geographic random baseline is constitutionally sound under federal law. The \textsc{bisect}
algorithm satisfies this standard, and the empirical evidence in I.1--I.3 confirms that
its incumbency outcomes are statistically indistinguishable from what neutral redistricting
would produce.

The state constitutional analysis reveals a spectrum of approaches to incumbency as
a redistricting criterion, ranging from explicit prohibition (Arizona, California
commission systems) to explicit permission (California commission guidelines, Ohio
pre-amendment criteria) to silence (most state constitutions). In all three categories,
the algorithmic indifference standard provides a defensible position: the algorithm
is neither prohibited by the explicit-prohibition states (because it neither protects
nor targets incumbents) nor required to protect incumbents by the explicit-permission
states (because permission is not mandate) nor constrained by the silent states (because
no constraint is stated).

For practitioners deploying \textsc{bisect} maps in redistricting litigation, the
central message is that the incumbency objection to algorithmic redistricting, while
rhetorically effective, is legally unfounded. No federal constitutional provision
requires incumbency protection; no VRA provision requires it; and no state constitutional
or statutory provision mandates it in any state we have studied. The algorithm's
indifference to incumbency is not a defect---it is a design choice that produces
incumbency outcomes equivalent to what a neutral geographic redistricting process would
produce, and those outcomes are constitutionally sound.

The most powerful response to the incumbency objection in litigation is therefore
empirical rather than doctrinal: the \textsc{bisect} map's incumbency outcomes are
within the range of geographic chance, while the enacted map's outcomes are significantly
below the baseline, demonstrating that enacted incumbent protection requires strategic
use of incumbency data that the algorithm does not access. This shifts the burden of
justification onto the enacted map---the party defending the map must explain why the
enacted map's near-zero pairing rate and above-baseline safe-seat fraction are a
geographic necessity rather than a strategic choice.
